\chapter{Risk Mathematics: Value at Risk and Expected Shortfall}\label{ch:risk}

\begin{objectives}
\begin{itemize}
\item define log-returns, Value at Risk and Expected Shortfall precisely, with sign conventions stated;
\item read the repository's historical-simulation estimator line by line and write it as order statistics;
\item place it among the Hyndman--Fan sample-quantile definitions and prove the inequalities between them;
\item work the ten-point example at $c\in\{0.95,0.90,0.80\}$ and the 503-return AAPL mock series;
\item state the coherence axioms, verify that VaR violates subadditivity, and prove that ES satisfies it;
\item explain why Basel's FRTB moved from 99\% VaR to 97.5\% ES.
\end{itemize}
{\small\color{muted}\textbf{Chapter references:} Artzner et al.~\cite{artzner1999coherent}; Acerbi and Tasche~\cite{acerbi2002coherence};
Rockafellar and Uryasev~\cite{rockafellar2000cvar,rockafellar2002general}; Hyndman and Fan~\cite{hyndman1996sample};
McNeil, Frey and Embrechts~\cite{mcneil2015qrm}; Jorion~\cite{jorion2007var}; Kupiec~\cite{kupiec1995techniques};
BCBS~\cite{bcbs1996amendment,bcbs2019frtb}.}
\end{objectives}

\section{Log-returns}
\begin{definition}[Log-return]\label[definition]{def:logret}
For positive prices $P_0,\dots,P_n$, the log-returns are $r_t=\ln(P_t/P_{t-1})$, $t=1,\dots,n$.
\end{definition}
Log-returns add over time: $\sum_{t=a+1}^{b}r_t=\ln(P_b/P_a)$. A log-return $r$ corresponds to the simple return
$e^r-1$, so a log-return of $-0.05$ is a loss of $1-e^{-0.05}=4.88\%$ of value. The repository computes VaR on
whatever series it is given; when that series is log-returns, the reported number is a loss in log units.

\repo[firstline=50,lastline=72]{src/statistics/risk.cpp}{\cc{risk::log_returns}. Every price must be positive (the
logarithm requires it), and $n+1$ prices yield $n$ returns.}{lst:logret}

\section{Value at Risk}
\begin{definition}[Value at Risk]\label[definition]{def:var}
Let $X$ be a return (profit positive, loss negative) with distribution function $F_X$, and let $c\in(0,1)$ be the
confidence level. The Value at Risk at level $c$ is
\[ \VaR_c(X) \;=\; -\inf\{x\in\R : F_X(x) > 1-c\}. \]
\end{definition}

\begin{layman}
Line up a thousand past trading days from worst to best. The 95\% VaR is roughly ``how bad the 50th-worst day was''.
It says nothing about how bad the 49 days beyond it were, and that blind spot is the subject of \cref{sec:es}.
\end{layman}

\begin{lemma}[Agreement with the loss convention]\label[lemma]{lem:lossconv}
With $L=-X$, \cref{def:var} equals $\inf\{\ell: \Prob(L>\ell)\le 1-c\}$, the loss-based definition of
McNeil, Frey and Embrechts~\cite{mcneil2015qrm}.
\end{lemma}
\begin{proof}
$\Prob(L>\ell)=\Prob(X<-\ell)$. Put $x=-\ell$: the condition $\Prob(X<x)\le 1-c$ holds exactly for $x\le q^+$, where
$q^+=\inf\{x:F_X(x)>1-c\}$ is the upper $(1-c)$-quantile (for $x<q^+$, $F_X(x)\le 1-c$; at $x=q^+$,
$\Prob(X<q^+)=\lim_{y\uparrow q^+}F_X(y)\le1-c$; for $x>q^+$, $\Prob(X<x)\ge F_X(y)>1-c$ for some $y\in(q^+,x)$). So the
set of admissible $\ell$ is $[-q^+,\infty)$ and its infimum is $-q^+=\VaR_c(X)$.
\end{proof}

\section{The repository's estimator}
\repo[firstline=13,lastline=48]{src/statistics/risk.cpp}{The shared helper. The tail size is
$k=\lfloor(1-c)\,n+10^{-9}\rfloor$; the epsilon is the subject of \cref{ch:float}. The comment states the post-condition:
the tail is \cc{sorted_returns[0, tail_index)}, its last element is the VaR order statistic, and its mean is CVaR.}{lst:tail}

\repo[firstline=102,lastline=111]{src/statistics/risk.cpp}{The two estimators. Both clamp at zero, so a sample whose worst
$k$ returns are all gains reports zero risk rather than a negative number.}{lst:varcvar}

\begin{definition}[Repository historical-simulation estimators]\label[definition]{def:repoest}
Given returns $X_1,\dots,X_n$ with order statistics $\ord1\le\dots\le\ord n$, confidence $c$, and
$k=\lfloor(1-c)n+\varepsilon\rfloor$ with $\varepsilon=10^{-9}$: if $k=0$ the estimators throw \cc{score::DomainError};
otherwise
\[ \widehat{\VaR}{}^{\,\mathrm{repo}}_c=\max\bigl(0,-\ord{k}\bigr),\qquad
   \widehat{\CVaR}{}^{\,\mathrm{repo}}_c=\max\Bigl(0,-\tfrac1k\textstyle\sum_{i=1}^{k}\ord{i}\Bigr). \]
\end{definition}

\section{Order statistics and the quantile zoo}
Plugging the empirical distribution $F_n$ into \cref{def:var} is the natural ``historical simulation'', but it is not what the
code computes, and it is worth knowing exactly how they differ. Hyndman and Fan catalogued nine sample-quantile
definitions in use in statistical packages~\cite{hyndman1996sample}; their \emph{type 1} is the inverse of the empirical
CDF. NumPy exposes it as \cc{method="inverted_cdf"}.

\begin{proposition}[Three empirical VaRs]\label[proposition]{prop:threevar}
Let $p=1-c$ and suppose $\varepsilon$ does not change $\lfloor pn\rfloor$ (\cref{thm:epsfloor}). Then, ignoring the clamp:
\begin{enumerate}[label=(\alph*)]
\item the plug-in of \cref{def:var} is $-\ord{\lfloor pn\rfloor+1}$;
\item the type-1 (lower) quantile $\inf\{x: F_n(x)\ge p\}$ gives $-\ord{\lceil pn\rceil}$;
\item the repository gives $-\ord{\lfloor pn\rfloor}$ (defined when $\lfloor pn\rfloor\ge1$).
\end{enumerate}
Hence $\widehat\VaR{}^{\mathrm{repo}}\ge\widehat\VaR{}^{\mathrm{type\,1}}\ge\widehat\VaR{}^{\mathrm{plug\text{-}in}}$, with
the first equality whenever $pn$ is an integer and the second whenever it is not.
\end{proposition}
\begin{proof}
$F_n(x)=j/n$ for $\ord j\le x<\ord{j+1}$ (ties only merge steps). (a) $F_n(x)>p\iff \#\{i:X_i\le x\}>pn\iff
\#\{\cdot\}\ge\lfloor pn\rfloor+1\iff x\ge\ord{\lfloor pn\rfloor+1}$, so the infimum is $\ord{\lfloor pn\rfloor+1}$.
(b) Likewise $F_n(x)\ge p\iff\#\{\cdot\}\ge\lceil pn\rceil$. (c) is \cref{def:repoest}. Since
$\lfloor pn\rfloor\le\lceil pn\rceil\le\lfloor pn\rfloor+1$ and order statistics are non-decreasing in their index,
negation reverses the inequalities. If $pn\in\N$ then $\lceil pn\rceil=\lfloor pn\rfloor$; otherwise
$\lceil pn\rceil=\lfloor pn\rfloor+1$.
\end{proof}

The repository is therefore \emph{conservative}: it reports the type-1 quantile at exact boundaries and one order statistic
further into the tail otherwise. For CVaR, the sample estimator consistent with Acerbi and Tasche~\cite{acerbi2002coherence}
averages the worst $pn$ observations, giving fractional weight to $\ord{\lfloor pn\rfloor+1}$:
\[ \widehat\ES_c=-\frac{1}{pn}\Bigl(\sum_{i=1}^{\lfloor pn\rfloor}\ord i+\bigl(pn-\lfloor pn\rfloor\bigr)\ord{\lfloor pn\rfloor+1}\Bigr), \]
which coincides with the repository's CVaR exactly when $pn$ is an integer.

\section{The ten-point worked example}\label{sec:tenpoint}
\begin{example}[Ten returns]\label[example]{ex:tenpoint}
Take $n=10$ returns
\[ X=\{0.00,\,0.01,\,0.02,\,0.03,\,0.04,\,-0.01,\,-0.02,\,-0.03,\,-0.04,\,-0.05\}. \]
Sorted ascending:
\[ \ord1,\dots,\ord{10}=-0.05,\,-0.04,\,-0.03,\,-0.02,\,-0.01,\,0.00,\,0.01,\,0.02,\,0.03,\,0.04. \]
\begin{itemize}
\item $c=0.95$: $pn$ evaluates to $0.50000000000000044$, so $k=0$ and both functions throw \cc{DomainError}: ten
observations cannot resolve a 5\% tail under this estimator. (The plug-in of \cref{def:var} would report $0.05$; the code
refuses instead.)
\item $c=0.90$: $k=1$. $\VaR=-\ord1=0.05$ and $\CVaR=-\ord1=0.05$. The tail is a single observation, so the two coincide.
\item $c=0.80$: $k=2$. $\VaR=-\ord2=0.04$ and $\CVaR=-(\ord1+\ord2)/2=-(-0.05-0.04)/2=0.045$.
\end{itemize}
\end{example}

\begin{progout}{risk\_check: the repository's own risk.cpp, compiled and called}{lst:riskout}
c=0.95 THROWS risk::historical_var: not enough returns (10) to resolve the 0.950000 confidence quantile
c=0.90 VaR=0.050000000000000003 CVaR=0.050000000000000003
c=0.85 VaR=0.050000000000000003 CVaR=0.050000000000000003
c=0.80 VaR=0.040000000000000001 CVaR=0.044999999999999998
c=0.75 VaR=0.040000000000000001 CVaR=0.044999999999999998
c=0.70 VaR=0.029999999999999999 CVaR=0.040000000000000001
c=0.50 VaR=0.01 CVaR=0.029999999999999999
\end{progout}

The 17-significant-digit output shows the binary64 values nearest $0.05$, $0.04$ and $0.045$; for example $0.05$ is
stored as $0.05000000000000000277\ldots$. \Cref{tab:quantiles} compares the estimators of \cref{prop:threevar} on the same data,
computed with exact rational arithmetic and with NumPy~2.5.1.

\begin{center}
\small
\begin{tabular}{@{}cccccccc@{}}
\toprule
\textcolor{muted}{$c$} & \textcolor{muted}{$k$} & \textcolor{muted}{repo VaR} & \textcolor{muted}{repo CVaR} & \textcolor{muted}{plug-in (a)} & \textcolor{muted}{type 1 (b)} & \textcolor{muted}{\texttt{inverted\_cdf}} & \textcolor{muted}{A--T ES}\\
\midrule
0.95 & 0 & throws & throws & 0.05 & 0.05 & 0.0500 & 0.0500\\
0.90 & 1 & 0.05 & 0.05 & 0.04 & 0.05 & 0.0500 & 0.0500\\
0.85 & 1 & 0.05 & 0.05 & 0.04 & 0.04 & 0.0400 & 0.0467\\
0.80 & 2 & 0.04 & 0.045 & 0.03 & 0.04 & 0.0400 & 0.0450\\
0.75 & 2 & 0.04 & 0.045 & 0.03 & 0.03 & 0.0300 & 0.0420\\
0.70 & 3 & 0.03 & 0.04 & 0.02 & 0.03 & 0.0200\rlap{$^\ast$} & 0.0400\\
0.50 & 5 & 0.01 & 0.03 & 0.00 & 0.01 & 0.0100 & 0.0300\\
\bottomrule
\end{tabular}
\captionof{table}{Estimators on the ten-point sample (\file{verify/var.py}). Rows with $pn\in\N$ ($c=0.90,0.80,0.70,0.50$)
show repo VaR $=$ type 1 and repo CVaR $=$ A--T ES, as \cref{prop:threevar} predicts. $^\ast$NumPy evaluates the
quantile at the binary64 value of $1-c$, here $0.30000000000000004>0.3$, and so steps to the next order statistic: the
epsilon problem of \cref{ch:float}, in a library.}\label{tab:quantiles}
\end{center}

\begin{figure}[htbp]
\centering
\begin{tikzpicture}
\begin{axis}[darkaxis,width=0.9\linewidth,height=6.3cm,xmin=-0.06,xmax=0.05,ymin=0,ymax=1.05,
  xlabel={return $x$},ylabel={$F_{10}(x)$},grid=major,xtick={-0.05,-0.04,-0.03,-0.02,-0.01,0,0.01,0.02,0.03,0.04},
  xticklabel style={/pgf/number format/fixed,/pgf/number format/precision=2},legend pos=south east]
\addplot[const plot,thick,acc] coordinates {(-0.06,0) (-0.05,0.1) (-0.04,0.2) (-0.03,0.3) (-0.02,0.4) (-0.01,0.5) (0,0.6) (0.01,0.7) (0.02,0.8) (0.03,0.9) (0.04,1.0) (0.05,1.0)};
\addlegendentry{empirical CDF}
\addplot[dashed,yel,thick] coordinates {(-0.06,0.1) (0.05,0.1)}; \addlegendentry{$p=0.10$}
\addplot[dashed,org,thick] coordinates {(-0.06,0.2) (0.05,0.2)}; \addlegendentry{$p=0.20$}
\addplot[only marks,mark=*,red,mark size=2.6pt] coordinates {(-0.05,0.1) (-0.04,0.2)};
\node[text=red,font=\footnotesize,anchor=south west,inner sep=1pt] at (axis cs:-0.037,0.47) {$\ord1$: VaR at $c=0.90$};
\draw[-{Stealth},red] (axis cs:-0.0375,0.49) -- (axis cs:-0.049,0.13);
\node[text=red,font=\footnotesize,anchor=south west,inner sep=1pt] at (axis cs:-0.037,0.68) {$\ord2$: VaR at $c=0.80$};
\draw[-{Stealth},red] (axis cs:-0.0375,0.70) -- (axis cs:-0.0398,0.23);
\end{axis}
\end{tikzpicture}
\caption{The empirical CDF of the ten returns (pgfplots). At $p=0.10$ the step reaches exactly $0.1$ at $\ord1$; the plug-in
of \cref{def:var} requires $F_n>p$ strictly and moves to $\ord2$, while type 1 and the repository stay at $\ord1$.}
\label{fig:ecdf}
\end{figure}

\section{A realistic series: the AAPL mock data}\label{sec:mockvar}
We generated the mock data in memory with \cc{gen_mock_data.generate()} (seed 42, 504 bars; no file was written), took the
AAPL closes, and ran the repository's \cc{log_returns}, \cc{calculate_risk_from_returns}, \cc{historical_var} and
\cc{historical_cvar} on the $503$ log-returns (\file{verify/mock_var.cpp}). The sample mean is $-0.000843$ and the sample
standard deviation $0.017694$.

\begin{center}
\small
\begin{tabular}{@{}ccccc@{}}
\toprule
\textcolor{muted}{$c$} & \textcolor{muted}{$k=\lfloor 503(1-c)+\varepsilon\rfloor$} & \textcolor{muted}{historical VaR} & \textcolor{muted}{historical CVaR} & \textcolor{muted}{Gaussian VaR $-(\hat\mu-z_c\hat\sigma)$}\\
\midrule
0.950 & 25 & 0.032002 & 0.038241 & 0.029948\\
0.975 & 12 & 0.036142 & 0.043080 & 0.035523\\
0.990 & 5 & 0.042478 & 0.050754 & 0.042006\\
\bottomrule
\end{tabular}
\captionof{table}{Repository estimators on the AAPL mock series, with the parametric normal VaR for comparison
($z_{0.95}=1.644854$, $z_{0.975}=1.959964$, $z_{0.99}=2.326348$).}\label{tab:mockvar}
\end{center}

The data are Gaussian by construction, so historical and parametric VaR should agree up to sampling noise, and they do
to within about 0.002. On real equity returns, whose tails are heavier than normal, the historical figure at 99\% is
typically the larger of the two.

\begin{figure}[htbp]
\centering
\begin{tikzpicture}
\begin{axis}[darkaxis,width=0.95\linewidth,height=6.2cm,ybar interval,xmin=-0.07,xmax=0.05,ymin=0,
  xlabel={daily log-return},ylabel={count (of 503)},xticklabel style={/pgf/number format/fixed,/pgf/number format/precision=2,/pgf/number format/fixed zerofill},
  xtick={-0.06,-0.04,-0.02,0,0.02,0.04}]
\addplot[fill=acc!45!bg,draw=acc] table[x=x,y=y]{data/aapl_hist.dat};
\draw[red,very thick] (axis cs:-0.032002,0) -- (axis cs:-0.032002,58) node[above,font=\footnotesize,text=red]{$-\VaR_{0.95}=-0.0320$};
\draw[org,very thick,dashed] (axis cs:-0.038241,0) -- (axis cs:-0.038241,44) node[above left,font=\footnotesize,text=org,xshift=4pt]{$-\CVaR_{0.95}=-0.0382$};
\end{axis}
\end{tikzpicture}
\caption{Histogram of the 503 AAPL mock log-returns in bins of 0.005 (pgfplots). The 95\% VaR cut leaves the worst 25
returns to its left; CVaR is their mean.}
\label{fig:aaplhist}
\end{figure}

\begin{figure}[htbp]
\centering
\begin{tikzpicture}
\begin{axis}[darkaxis,width=0.9\linewidth,height=5.8cm,xlabel={confidence $c$},ylabel={loss (log-return units)},
  grid=major,legend pos=north west,xticklabel style={/pgf/number format/fixed,/pgf/number format/precision=3}]
\addplot[const plot,thick,red] table[x=c,y=var]{data/aapl_varcurve.dat}; \addlegendentry{historical VaR}
\addplot[const plot,thick,org] table[x=c,y=cvar]{data/aapl_varcurve.dat}; \addlegendentry{historical CVaR}
\end{axis}
\end{tikzpicture}
\caption{VaR and CVaR of the AAPL mock series as functions of $c\in[0.90,0.995]$ in steps of $0.0025$ (pgfplots). Both are
step functions because $k$ is an integer; CVaR always lies above VaR.}
\label{fig:varcurve}
\end{figure}

\section{Coherent risk measures}\label{sec:es}
What should a risk measure satisfy? Artzner, Delbaen, Eber and Heath answered with four axioms~\cite{artzner1999coherent}.
Write losses as random variables $L$ (positive = bad) and a risk measure as $\rho(L)\in\R$, the capital needed to make
the position acceptable.

\begin{definition}[Coherent risk measure]\label[definition]{def:coherent}
$\rho$ is \emph{coherent} if for all losses $L,L_1,L_2$, constants $m\in\R$ and $\lambda\ge0$:
\begin{enumerate}[label=\textbf{(C\arabic*)}]
\item \emph{Monotonicity}: $L_1\le L_2$ almost surely implies $\rho(L_1)\le\rho(L_2)$.
\item \emph{Translation invariance}: $\rho(L+m)=\rho(L)+m$.
\item \emph{Positive homogeneity}: $\rho(\lambda L)=\lambda\rho(L)$.
\item \emph{Subadditivity}: $\rho(L_1+L_2)\le\rho(L_1)+\rho(L_2)$.
\end{enumerate}
\end{definition}

Subadditivity is the axiom of diversification: merging two books should never \emph{create} risk. A desk-level limit
system built on a subadditive measure can add up desk risks and know the total is an upper bound.

\begin{theorem}[VaR is not subadditive]\label{thm:varnotsub}
There exist losses $L_1,L_2$ with $\VaR_{0.95}(L_1+L_2)>\VaR_{0.95}(L_1)+\VaR_{0.95}(L_2)$.
\end{theorem}
\begin{proof}
Let $L_1,L_2$ be independent, each equal to $100$ with probability $0.04$ (a loan defaults) and $0$ otherwise. Using the
loss form of \cref{lem:lossconv}, $\VaR_{0.95}(L)=\inf\{\ell:\Prob(L>\ell)\le0.05\}$. For each loan,
$\Prob(L_i>0)=0.04\le0.05$, so $\VaR_{0.95}(L_i)=0$. For the sum,
\begin{gather*} \Prob(L_1+L_2=0)=0.96^2=\tfrac{576}{625},\qquad \Prob(L_1+L_2=100)=2(0.04)(0.96)=\tfrac{48}{625},\\ \Prob(L_1+L_2=200)=0.04^2=\tfrac{1}{625}. \end{gather*}
Then $\Prob(L_1+L_2>\ell)=\frac{49}{625}=0.0784>0.05$ for $0\le\ell<100$, and $=\frac1{625}=0.0016\le0.05$ for
$100\le\ell<200$. So $\VaR_{0.95}(L_1+L_2)=100>0+0$.
\end{proof}

\begin{verified}
\file{verify/var.py} recomputes the three probabilities with exact fractions ($576/625$, $48/625$, $1/625$). A Monte
Carlo check with $2\times10^6$ draws per loan (seed 7) gives empirical 95\% quantiles of $0$, $0$ and $100$, with default
frequencies $0.0399$ and $0.0400$ and $\Prob(\text{sum}>0)=0.0783$.
\end{verified}

The pathology is not exotic: any portfolio of rare, large, independent losses (credit, operational risk, short
out-of-the-money options) can hide below a VaR threshold individually and breach it together. VaR \emph{is} subadditive
for jointly elliptical distributions such as the multivariate normal~\cite{mcneil2015qrm}, which is part of why the problem
went unnoticed in Gaussian models.

\begin{definition}[Expected Shortfall]\label[definition]{def:es}
For a loss $L$ with $\E|L|<\infty$ and $\alpha\in(0,1)$,
$\ES_\alpha(L)=\frac{1}{1-\alpha}\int_\alpha^1 \VaR_u(L)\,du$, the average of the VaRs beyond level $\alpha$.
For continuous $L$ this equals $\E[L\mid L\ge\VaR_\alpha(L)]$, the ``conditional VaR'' of the repository's naming.
\end{definition}

\begin{theorem}[Rockafellar--Uryasev representation~\cite{rockafellar2000cvar,rockafellar2002general}]\label{thm:ru}
For $L$ with $\E|L|<\infty$,
\[ \ES_\alpha(L)=\min_{t\in\R}\Bigl\{\,t+\frac{1}{1-\alpha}\E\bigl[(L-t)^+\bigr]\Bigr\}, \]
and the minimum is attained at $t=\VaR_\alpha(L)$.
\end{theorem}

\begin{theorem}[ES is coherent~\cite{acerbi2002coherence}]\label{thm:escoherent}
$\ES_\alpha$ satisfies (C1)--(C4).
\end{theorem}
\begin{proof}
Write $g_L(t)=t+\frac1{1-\alpha}\E[(L-t)^+]$, so $\ES_\alpha(L)=\min_t g_L(t)$ by \cref{thm:ru}.
(C1) If $L_1\le L_2$ a.s.\ then $(L_1-t)^+\le(L_2-t)^+$, so $g_{L_1}\le g_{L_2}$ pointwise, and minima preserve the order.
(C2) $g_{L+m}(t+m)=t+m+\frac1{1-\alpha}\E[(L-t)^+]=g_L(t)+m$; minimising over $t$ gives $\ES(L+m)=\ES(L)+m$.
(C3) For $\lambda>0$, $g_{\lambda L}(\lambda t)=\lambda g_L(t)$; for $\lambda=0$ both sides are $0$.
(C4) For any $t_1,t_2$, since $(a+b)^+\le a^++b^+$,
\[ \ES_\alpha(L_1+L_2)\le g_{L_1+L_2}(t_1+t_2)\le t_1+t_2+\tfrac1{1-\alpha}\bigl(\E(L_1-t_1)^++\E(L_2-t_2)^+\bigr)=g_{L_1}(t_1)+g_{L_2}(t_2). \]
Taking $t_i=\VaR_\alpha(L_i)$, the minimisers, gives $\ES_\alpha(L_1+L_2)\le\ES_\alpha(L_1)+\ES_\alpha(L_2)$.
\end{proof}

\begin{example}[The loans again, with ES]\label[example]{ex:loanses}
At $\alpha=0.95$ the worst 5\% of probability mass of one loan is $0.04$ at $100$ and $0.01$ at $0$, so
$\ES_{0.95}(L_i)=(0.04\cdot100)/0.05=80$. For the sum, the worst 5\% is $0.0016$ at $200$ and $0.0484$ at $100$:
$\ES_{0.95}(L_1+L_2)=(0.0016\cdot200+0.0484\cdot100)/0.05=103.2\le160$. Subadditivity holds, and ES, unlike VaR, sees the
default risk of a single loan. (Exact fractions in \file{verify/var.py}: $80$ and $516/5$.)
\end{example}

\section{From Basel to FRTB}
The 1996 amendment to the Basel Capital Accord let banks compute market-risk capital from internal 99\%, 10-day VaR
models~\cite{bcbs1996amendment}. After 2008 the Basel Committee's Fundamental Review of the Trading Book replaced
99\% VaR with 97.5\% Expected Shortfall~\cite{bcbs2019frtb}. The level was chosen so that for normal returns the two are
close: with $\varphi$ the standard normal density,
\[ \ES_{0.975}=\frac{\varphi(z_{0.975})}{0.025}=2.337803 \qquad\text{vs.}\qquad \VaR_{0.99}=z_{0.99}=2.326348 \]
standard deviations (computed with Python's \cc{statistics.NormalDist}). Under normality the capital barely moves;
under fat tails ES grows while VaR does not look. Jorion's textbook~\cite{jorion2007var} treats the VaR side of this
history in depth; \cref{sec:hist-ltcm} recounts the episodes that discredited VaR as a sole limit.

\begin{exercises}
\ex{ex:comment}{\lvA} An earlier version of the comment above \cc{struct HistoricalTail} (\cref{lst:tail}) claimed that
``the index-0 element of the tail is the VaR quantile''; it was corrected while this tutorial was being written, and the
code never changed. Find $n$ and $c$ for which the two readings, $-\ord1$ versus $-\ord k$, give different VaRs, and
state for which $k$ they always agree.

\ex{ex:atcvar}{\lvB} Implement the Acerbi--Tasche sample ES (the displayed formula after \cref{prop:threevar}) with the same
exception conventions as \cc{historical_cvar}. Check it against \cref{tab:quantiles} and decide whether it should
replace the current CVaR, or live beside it.

\ex{ex:clamp}{\lvB} Show that the clamp $\max(0,\cdot)$ breaks translation invariance (C2) of the estimator. Give a
pre-trade use case where a \emph{negative} VaR (a position that is profitable even in its 5\% worst case) is meaningful,
and propose an API that serves both users.

\ex{ex:kupiec}{\lvC} Kupiec's proportion-of-failures test~\cite{kupiec1995techniques} checks that VaR exceedances occur
with frequency $1-c$. Count in-sample exceedances of the 95\% VaR in the AAPL mock series and explain why the count is
fixed by construction. Design a rolling out-of-sample backtest using \cc{ColumnStore::scan}, with a 250-day window.

\ex{ex:ellip}{\lvB} Prove that if $(L_1,L_2)$ is bivariate normal with mean zero, then
$\VaR_\alpha(L_1+L_2)\le\VaR_\alpha(L_1)+\VaR_\alpha(L_2)$ for $\alpha\ge\frac12$.

\ex{ex:sqrt}{\lvB} Basel's 10-day VaR was often computed as $\sqrt{10}$ times the 1-day figure. Under what assumptions on
daily log-returns is this exact? Which assumption does volatility clustering break, and in which direction is the error?
\end{exercises}

